\section{Discussion}
\label{sec:discussion}

This work trained Barlow Twins on nested Tiny-ImageNet subsets under fixed
generous budgets with downstream-blind checkpoint selection, and measured
each selected encoder by a CIFAR-10 linear probe against a
random-initialization floor. Three observations stand out: pre-training
pays at every set size (19.9 to 30.0 points over the floor), the curve dips
at 4k--8k, and the dip coincides with the largest measured divergence
between the selection signal and downstream accuracy.

\paragraph{What the dip means.}
The naive reading of a data-efficiency curve is that each point reflects
the value of that much data. Our results say the points also reflect the
\emph{stopping and selection policy}: where the in-domain signal and the
transfer optimum diverge, a downstream-blind protocol ships a checkpoint
well past the transfer peak, and at the diagnostic level the 4k--16k cost
of doing so (up to 4.7 points) exceeds what the added data buys. The
matched-compute slice of Section~\ref{sec:results_decoupling} is the
sharpest evidence that the data itself is not at fault: read at an equal
step count, before the protocols diverge, more data monotonically helps. The practical reading, stated
tentatively: in this regime, what a practitioner measures to stop and
select is as consequential as how much unlabeled data they have. We would
not have seen this with downstream-based selection --- selecting on the
downstream signal hides the divergence by construction, at the price of
tuning the shipped checkpoint to the evaluation task; the
downstream-selected campaign in Section~\ref{sec:results_decoupling} gives
a budget-confounded estimate of that tuning from the other side.

\paragraph{Does the family-level robustness ranking apply here?}
The report that joint-embedding methods need more pre-training data than
masked-prediction methods rests on fine-tuning
evaluation~\cite{elnouby2021largescale}, and the same study finds denoising
autoencoders comparatively weak under linear probing. Whether the ranking
carries over to linear-probe evaluation at our scale is therefore open. Our
curve contributes one joint-embedding data point measured under a linear
probe, at set sizes one to two orders of magnitude below those tested
there --- with the added nuance that at small scale, the measured
``robustness to set size'' depends on the selection policy as much as on
the method. The five-method group curve this study is part of includes a
masked-autoencoder arm under the same protocol and will offer a first
side-by-side look, with the caveat that per-method tuning differs across
students.

\paragraph{Why augmentations stayed fixed.}
If augmentations act chiefly as artificial dataset
enlargement~\cite{moutakanni2024you}, strengthening them is the obvious
lever exactly where data is scarce, and it is fair to ask why we left them
alone. An augmentation is a commitment: a joint-embedding method learns to
be invariant to whatever its augmentations vary~\cite{xiao2021what}, so a
stronger augmentation buys small-$N$ robustness by discarding the
distinctions it augments away --- a representation made invariant to color
cannot separate red from yellow cars~\cite{xiao2021what} --- and which
distinctions a downstream task needs is unknown at pre-training time in a
pre-train-once, probe-anywhere setting. Augmentation strength also
interacts with $N$ itself, so varying both would confound the curve.
Varying a capacity knob instead --- the Barlow Twins projector width ---
leaves data, views, and objective untouched; that study is scoped as the
immediate next step (Section~\ref{sec:conclusion}).

\paragraph{Compute.}
The eight stable phases of the campaign total 10.2 GPU-hours on one
consumer GPU; each cooldown adds minutes (up to roughly 25 for the largest
split, sized as 20\% of its selected step), and each linear probe takes
about 6 minutes. We report these numbers so that others can budget similar
studies.
